\chapter{Requerimientos del Payload}
\label{ch:requerimientos}

Este capítulo presenta los requerimientos formales del
\textbf{CMOS Microscopy Payload (CMOSM)}, siguiendo lineamientos de ingeniería de
sistemas ECSS/NASA y alineado con el \textit{System Requirements Document} (SRD)
del proyecto INTISAT. Los requerimientos se clasifican en funcionales, de
rendimiento, ópticos, eléctricos, ambientales y operacionales, e incluyen su
método de verificación asociado.

% -------------------------------------------------------------------
\section{Clasificación de métodos de verificación}

Los métodos de verificación utilizados en este documento se definen como:

\begin{itemize}
    \item \textbf{T} -- Test: ensayo directo de hardware y/o software.
    \item \textbf{A} -- Analysis: análisis, simulación o evaluación numérica.
    \item \textbf{I} -- Inspection: inspección visual o revisión física.
    \item \textbf{R} -- Review of Design: revisión documental o de código.
\end{itemize}

% -------------------------------------------------------------------
\section{Requerimientos funcionales (FR)}

\setlength{\LTpre}{0.5\baselineskip}
\setlength{\LTpost}{0.5\baselineskip}
\renewcommand{\arraystretch}{1.3}
\begin{longtable}{p{1.2cm} p{10.3cm} p{2.5cm}}
\toprule
\textbf{ID} & \textbf{Descripción} & \textbf{Verificación} \\
\midrule
FR-01 & El payload debe capturar imágenes a resolución nativa de
2592$\times$1944 píxeles utilizando un sensor CMOS OV5647. & T, R \\
FR-02 & El sistema debe permitir el control de iluminación mediante patrones
programables mostrados en una pantalla OLED. & T, I \\
FR-03 & El sistema debe generar un archivo de metadatos por cada captura,
incluyendo parámetros de adquisición y marca temporal. & T, R \\
FR-04 & El sistema debe permitir operación en modo manual y semiautomático. & T \\
FR-05 & El payload debe integrarse con una Raspberry Pi actuando como
computadora de carga útil y permitir interacción funcional con el sistema OBDH
del proyecto INTISAT. & I, R \\
FR-06 & El software debe permitir modificar parámetros de la iluminación OLED,
incluyendo posición, tamaño e intensidad del patrón. & T \\
\bottomrule
\end{longtable}
\renewcommand{\arraystretch}{1}

% -------------------------------------------------------------------
\section{Requerimientos de rendimiento (PR)}
\renewcommand{\arraystretch}{1.3}

\begin{longtable}{p{1.2cm} p{10.3cm} p{2.5cm}}
\toprule
\textbf{ID} & \textbf{Descripción} & \textbf{Verificación} \\
\midrule
PR-01 & El sistema debe generar imágenes con contraste suficiente para permitir
la observación visual de estructuras celulares planas bajo condiciones nominales
de laboratorio. & T, I \\
PR-02 & La imagen base capturada debe presentar un nivel de ruido compatible con
análisis visual, evaluado mediante métricas estadísticas básicas. & A, T \\
PR-03 & El histograma de la imagen capturada no debe presentar más del 1\% de
píxeles saturados bajo condiciones nominales de exposición. & T \\
PR-04 & La variación relativa de intensidad luminosa entre capturas consecutivas
debe ser inferior al 5\% bajo condiciones estables. & T \\
PR-05 & Las técnicas computacionales experimentales evaluadas no deben degradar
la información presente en la imagen base. & A, I \\
\bottomrule
\end{longtable}
\renewcommand{\arraystretch}{1}

% -------------------------------------------------------------------
\section{Requerimientos ópticos (OR)}
\renewcommand{\arraystretch}{1.3}

\begin{longtable}{p{1.2cm} p{10.3cm} p{2.5cm}}
\toprule
\textbf{ID} & \textbf{Descripción} & \textbf{Verificación} \\
\midrule
OR-01 & El sistema debe operar en configuración de microscopía sin lentes, sin
emplear elementos ópticos refractivos para la formación de imagen. & R \\
OR-02 & El diseño debe permitir colocar la muestra en contacto cercano con el
plano del sensor CMOS. & I \\
OR-03 & El sistema de iluminación debe proporcionar una iluminación
aproximadamente uniforme sobre el área de interés. & T \\
OR-04 & La arquitectura óptica no debe requerir mecanismos de enfoque activo. & R \\
OR-05 & La iluminación estructurada se implementa únicamente como modo
experimental y no afecta la operación nominal. & R \\
\bottomrule
\end{longtable}
\renewcommand{\arraystretch}{1}

% -------------------------------------------------------------------
\section{Requerimientos eléctricos (ER)}
\renewcommand{\arraystretch}{1.3}

\begin{longtable}{p{1.2cm} p{10.3cm} p{2.5cm}}
\toprule
\textbf{ID} & \textbf{Descripción} & \textbf{Verificación} \\
\midrule
ER-01 & El consumo total del payload debe ser inferior a 1.5 W en operación
nominal. & A, T \\
ER-02 & El sistema de iluminación OLED debe operar mediante bus I2C a nivel
lógico de 3.3 V. & I \\
ER-03 & La Raspberry Pi debe proveer alimentación y control al sensor CMOS. & I, R \\
ER-04 & El diseño debe incluir protecciones básicas contra descargas
electrostáticas en conexiones externas. & I \\
\bottomrule
\end{longtable}
\renewcommand{\arraystretch}{1}

% -------------------------------------------------------------------
\section{Requerimientos ambientales (ENR)}
\renewcommand{\arraystretch}{1.3}

\begin{longtable}{p{1.2cm} p{10.3cm} p{2.5cm}}
\toprule
\textbf{ID} & \textbf{Descripción} & \textbf{Verificación} \\
\midrule
ENR-01 & El sistema debe operar en el rango de temperatura de 10--40~$^\circ$C en
condiciones de laboratorio. & T \\
ENR-02 & Los componentes deben resistir vibraciones de laboratorio de hasta
3.5~g RMS. & A, T \\
ENR-03 & El sensor CMOS debe contar con protección mecánica básica frente a polvo. & I \\
ENR-04 & El sistema de iluminación debe mantener contraste operativo hasta
40~$^\circ$C. & T \\
\bottomrule
\end{longtable}
\renewcommand{\arraystretch}{1}

% -------------------------------------------------------------------
\section{Requerimientos operacionales (OPR)}
\renewcommand{\arraystretch}{1.3}

\begin{longtable}{p{1.2cm} p{10.3cm} p{2.5cm}}
\toprule
\textbf{ID} & \textbf{Descripción} & \textbf{Verificación} \\
\midrule
OPR-01 & El sistema debe iniciar su operación en menos de 15 s desde el
encendido. & T \\
OPR-02 & El usuario debe poder modificar parámetros de iluminación en tiempo
cuasi--real ($\leq$ 100 ms). & T \\
OPR-03 & El sistema debe permitir la ejecución de modos experimentales de captura
sin interferir con el modo nominal. & R \\
OPR-04 & El sistema debe disponer de un modo seguro de apagado. & R \\
OPR-05 & El sistema debe registrar logs de operación, incluyendo errores,
parámetros y eventos relevantes. & R \\
\bottomrule
\end{longtable}
\renewcommand{\arraystretch}{1}

% -------------------------------------------------------------------
\section{Resumen}

El conjunto de requerimientos definidos en este capítulo establece el marco de
desempeño esperado para el CMOS Microscopy Payload en su fase actual de
desarrollo. Cada requerimiento será verificado mediante los métodos descritos en
el Capítulo~\ref{ch:verificacion}, asegurando coherencia entre diseño,
implementación y validación experimental.
